\documentclass{amsart}
% For dealing with references we use the comment environment
\usepackage{verbatim}
\newenvironment{reference}{\comment}{\endcomment}
%\newenvironment{reference}{}{}
\newenvironment{slogan}{\comment}{\endcomment}
\newenvironment{history}{\comment}{\endcomment}

% For commutative diagrams we use Xy-pic
\usepackage[all]{xy}

% We use 2cell for 2-commutative diagrams.
\xyoption{2cell}
\UseAllTwocells

% We use multicol for the list of chapters between chapters
\usepackage{multicol}

% This is generally recommended for better output
\usepackage{lmodern}
\usepackage[T1]{fontenc}

% For cross-file-references

% Package for hypertext links:
\usepackage{hyperref}

% For any local file, say "hello.tex" you want to link to please

% Theorem environments.
%
\theoremstyle{plain}
\newtheorem{theorem}[subsection]{Theorem}
\newtheorem{proposition}[subsection]{Proposition}
\newtheorem{lemma}[subsection]{Lemma}

\theoremstyle{definition}
\newtheorem{definition}[subsection]{Definition}
\newtheorem{example}[subsection]{Example}
\newtheorem{exercise}[subsection]{Exercise}
\newtheorem{situation}[subsection]{Situation}

\theoremstyle{remark}
\newtheorem{remark}[subsection]{Remark}
\newtheorem{remarks}[subsection]{Remarks}

\numberwithin{equation}{subsection}

% Macros
%
\def\lim{\mathop{\mathrm{lim}}\nolimits}
\def\colim{\mathop{\mathrm{colim}}\nolimits}
\def\Spec{\mathop{\mathrm{Spec}}}
\def\Hom{\mathop{\mathrm{Hom}}\nolimits}
\def\Ext{\mathop{\mathrm{Ext}}\nolimits}
\def\SheafHom{\mathop{\mathcal{H}\!\mathit{om}}\nolimits}
\def\SheafExt{\mathop{\mathcal{E}\!\mathit{xt}}\nolimits}
\def\Sch{\mathit{Sch}}
\def\Mor{\mathop{\mathrm{Mor}}\nolimits}
\def\Ob{\mathop{\mathrm{Ob}}\nolimits}
\def\Sh{\mathop{\mathit{Sh}}\nolimits}
\def\NL{\mathop{N\!L}\nolimits}
\def\CH{\mathop{\mathrm{CH}}\nolimits}
\def\proetale{{pro\text{-}\acute{e}tale}}
\def\etale{{\acute{e}tale}}
\def\QCoh{\mathit{QCoh}}
\def\Ker{\mathop{\mathrm{Ker}}}
\def\Im{\mathop{\mathrm{Im}}}
\def\Coker{\mathop{\mathrm{Coker}}}
\def\Coim{\mathop{\mathrm{Coim}}}

% Boxtimes
%
\DeclareMathSymbol{\boxtimes}{\mathbin}{AMSa}{"02}

%
% Macros for moduli stacks/spaces
%
\def\QCohstack{\mathcal{QC}\!\mathit{oh}}
\def\Cohstack{\mathcal{C}\!\mathit{oh}}
\def\Spacesstack{\mathcal{S}\!\mathit{paces}}
\def\Quotfunctor{\mathrm{Quot}}
\def\Hilbfunctor{\mathrm{Hilb}}
\def\Curvesstack{\mathcal{C}\!\mathit{urves}}
\def\Polarizedstack{\mathcal{P}\!\mathit{olarized}}
\def\Complexesstack{\mathcal{C}\!\mathit{omplexes}}
% \Pic is the operator that assigns to X its picard group, usage \Pic(X)
% \Picardstack_{X/B} denotes the Picard stack of X over B
% \Picardfunctor_{X/B} denotes the Picard functor of X over B
\def\Pic{\mathop{\mathrm{Pic}}\nolimits}
\def\Picardstack{\mathcal{P}\!\mathit{ic}}
\def\Picardfunctor{\mathrm{Pic}}
\def\Deformationcategory{\mathcal{D}\!\mathit{ef}}

\setlength{\emergencystretch}{3em}
\begin{document}
\title[Stacks Project, Tag 0C2C]{417 --- Disconnected punctured spectra of completions and henselizations}
\maketitle
\noindent Copyright (C) 2005--2025 Johan de Jong. Stacks Project authors. Source: \href{https://stacks.math.columbia.edu/tag/0C2C}{Tag 0C2C}.

\noindent Editorial difficulty note (not author text). Difficulty H2: The proof encodes a disconnected completed punctured spectrum by two ideals with nilpotent product and finite-length conormal data. It algebraizes a quotient, uses henselian etale localization to descend its ideal, and reconstructs the disconnected decomposition before completion..

\noindent Editorial provenance note (not author text). History: excerpt from revision \texttt{a0444\allowbreak 6e57e\allowbreak c1fbc\allowbreak 252a8\allowbreak 71afc\allowbreak ec775\allowbreak 2fb28\allowbreak 07b14}, more-algebra.tex, lines 32144--32207. Reference presentation was adapted; original-author context and core support blocks were retained or added. The mathematical wording is unchanged. Licensed under GNU FDL 1.2 or later; no invariant sections or cover texts. See ../licenses/COPYING. Context blocks reproduce author definitions and supporting results; they are not additional counted problems.

\begin{lemma}
\label{lemma-completion-disconnected}
Let $(A, \mathfrak m)$ be a Noetherian local ring. The punctured spectrum
of $A^\wedge$ is disconnected if and only if the punctured spectrum of $A^h$
is disconnected.
\end{lemma}

\begin{proof}
Since the completion of $A$ and $A^h$ are the same, we may assume
that $A$ is henselian (Lemma \href{https://stacks.math.columbia.edu/tag/06LJ}{lemma henselization noetherian (Tag 06LJ)}).

\medskip\noindent
Since $A \to A^\wedge$ is faithfully flat (see reference just given)
the map from the punctured spectrum of $A^\wedge$ to the punctured
spectrum of $A$ is surjective
(see Algebra, Lemma \href{https://stacks.math.columbia.edu/tag/00HQ}{lemma ff rings (Tag 00HQ)}).
Hence if the punctured spectrum of $A$ is disconnected, then
the same is true for $A^\wedge$.

\medskip\noindent
Assume the punctured spectrum of $A^\wedge$ is disconnected.
This means that
$$
\Spec(A^\wedge) \setminus \{\mathfrak m^\wedge\} = Z \amalg Z'
$$
with $Z$ and $Z'$ closed. Let
$\overline{Z}, \overline{Z}' \subset \Spec(A^\wedge)$ be the closures.
Say $\overline{Z} = V(J)$, $\overline{Z}' = V(J')$ for some ideals
$J, J' \subset A^\wedge$. Then $V(J + J') = \{\mathfrak m^\wedge\}$
and $V(JJ') = \Spec(A^\wedge)$. The first equality means that
$\mathfrak m^\wedge = \sqrt{J + J'}$ which implies
$(\mathfrak m^\wedge)^e \subset J + J'$ for some $e \geq 1$.
The second equality implies every element
of $JJ'$ is nilpotent hence $(JJ')^n = 0$ for some $n \geq 1$.
Combined this means that $J^n/J^{2n}$ is annihilated by
$J^n$ and $(J')^n$ and hence by $(\mathfrak m^\wedge)^{2en}$.
Thus we may apply Lemma \href{https://stacks.math.columbia.edu/tag/0ALR}{lemma quotient by idempotent (Tag 0ALR)}
to see that there is a finite type $A$-algebra $C$ and an
isomorphism $A^\wedge/J^n = C^\wedge$.

\medskip\noindent
The rest of the proof is exactly the same as the second part
of the proof of Lemma \href{https://stacks.math.columbia.edu/tag/0C2B}{lemma one dimensional formal branch (Tag 0C2B)};
of course that lemma is a special case of this one!
We have $C/\mathfrak m C = A/\mathfrak m$ because $A^\wedge/J^n$ has the same
property. Hence $\mathfrak m_C = \mathfrak m C$ is a maximal ideal
and $A \to C$ is unramified at $\mathfrak m_C$
(Algebra, Lemma \href{https://stacks.math.columbia.edu/tag/02FM}{lemma characterize unramified (Tag 02FM)}).
After replacing $C$ by a principal localization we may
assume that $C$ is a quotient of an \'etale $A$-algebra $B$, see
Algebra, Proposition \href{https://stacks.math.columbia.edu/tag/0395}{proposition unramified locally standard (Tag 0395)}.
However, since the residue field extension of $A \to C_{\mathfrak m_C}$
is trivial and $A$ is henselian, we conclude that $B = A$
again after a localization.
Thus $C = A/I$ for some ideal $I \subset A$ and it follows that
$J^n = IA^\wedge$ (because completion is exact in our situation by
Algebra, Lemma \href{https://stacks.math.columbia.edu/tag/00MB}{lemma completion flat (Tag 00MB)}) and $I = J^n \cap A$
(by flatness of $A \to A^\wedge$).
By symmetry $I' = (J')^n \cap A$ satisfies $(J')^n = I'A^\wedge$.
Then $\mathfrak m^e \subset I + I'$ and $II' = 0$
and we conclude that $V(I)$ and $V(I')$ are closed subschemes
which give the desired disjoint union decomposition of the
punctured spectrum of $A$.
\end{proof}
\end{document}
